\chapter{Contextualizing the Napatan period: historical background and cultural interactions}


In this chapter, I will examine the historical context of the Napatan period, as well as the history of interactions between Nubia and Egypt. I begin with an outline of the period, tracing its origins from the end of the Egyptian occupation during the New Kingdom to the last known rulers of the Napatan era. Unfortunately, both the beginning and end of the period are characterized by a scarcity of material evidence. Notably, for the latter part of the period, we know little beyond the names of its rulers, leaving much about the region's overall history still unknown.


Within this context, it is necessary to explore the origins and early chronology of the Napatan dynasty---a question in the center of scholarly debate for many years. The various theories surrounding the topic are important in order to understand the processes of cultural entanglement in the Nile Valley region. Specifically, these discussions focus on whether the emergence of the Napatan rulers was an internal development or the result of a new wave of Egyptian immigrants.


Following this discussion and to set the stage for a detailed examination of amulets during the Napatan period, I will outline the complex and historical relations between Nubia and Egypt, which are particularly important for grasping why the concept of objectscapes offers a more nuanced framework for interpreting not only the presence of "Egyptian-style" features and material culture in Nubia but also the persistence of Nubian objects and motifs. Moreover, by focusing on recent archaeological discoveries and material evidence, we can gain insights into the experiences of non-elite social groups that have often been overlooked in traditional written and visual sources.


Amulets were also part of this complex interplay, and by examining the prolonged cultural interactions between Nubia and Egypt over many centuries and across diverse cultures, we will be better equipped to analyze and discuss the use of amulets in Napatan tombs across different social groups later on.


\section{Nubia during the Napatan Period}


The Napatan period can be characterized as a shift in Kush's position in the Nile Valley during the first half of the first millennium BCE. Following near five centuries of Egyptian rule (ca. 1550---1070 BCE) and two subsequent centuries with no clear archaeological record, the 8th century saw the rise of new Kushite rulers that conquered Upper Egypt and reunified the land as the 25th Dynasty (ca. 715---664 BCE).


The period is named after the burial sites of its kings, El-Kurru and Nuri, located near the ancient city of Napata \parencite[p.~1]{popeNapatanPeriod2020}. The origin and precise chronology of the Napatan period has been the subject of different theories and debates for many years. After the Egyptian withdrawal from Lower Nubia ca. 1070 BCE, material and textual evidence is extremely scarce in the region for the next three centuries. For many years, scholars have referred to the period following the New Kingdom as the "dark ages" of Nubia, due to the scarcity of archaeological evidence \parencites[p.~250]{adamsNubiaCorridorAfrica1977}[pp.~88, 109-110]{törökKingdomKushHandbook1997}[p.~129]{morkotBlackPharaohsEgypts2000}[pp.~412-413]{williamsNapatanNeoKushiteState2020}. Earlier authors even proposed that the temple towns of Lower Nubia were depopulated since the Egyptians had returned to their native land while Nubians withdrew to southern regions of Upper Nubia \parencites[p.~108-109]{adamsPostPharaonicNubiaLight1964}{firthArchaeologicalReportNubia1912}[p.109]{edwardsNubianArchaeologySudan2004}. Recently, it has been suggested that this was likely a misinterpretation of the archaeological material, which was being identified as either New Kingdom or Napatan \parencites[p.~37]{williamsExcavationsAbuSimbel1990}[p.~109-110]{törökKingdomKushHandbook1997}[p.~415]{williamsNapatanNeoKushiteState2020}. Indeed, recent archaeological work indicate that some cemeteries were occupied continuously from the New Kingdom to the Napatan period, such as Tombos \parencites{buzonEntanglementFormationAncient2016}[p.~7]{smithEthnicityCulture2007}, Hillat el-Arab \parencite{vincentelliHillatElArabJoint2006}, Qustul, Adindan \parencite[p.141]{williamsExcavationsAbuSimbel1992}, Debeira 176 \parencite[p.~200-205]{säve-söderberghScandinavianJointExpedition1989}, Amara West \parencites[p.~57]{spencerCemeteriesLateRamesside2009}{binderCemeteryAmaraWest2011}{binder10th9thCenturyBC2011}, Sanam \parencite{lohwasserAspekteNapatanischenGesellschaft2012}, and Attab and Ferka regions \parencite{budkaEndEgyptianNew2025}. Additionally, settlements and fortresses have also been dated to the post New Kingdom period \parencites[p.~1071-1078]{jesseFendingDesertDwellers2019}[p.~416-417]{williamsNapatanNeoKushiteState2020}.


Despite the limited evidence available, Lower Nubia likely remained populated after the Egyptian withdrawal. Given their role as administrators under Egyptian rule, it is plausible that local elites continued to govern independently and contributed to the rise of the Napatan state \parencites[p.~116]{edwardsNubianArchaeologySudan2004}[p.~111]{törökKingdomKushHandbook1997}. An inscription discovered at the Egyptian fortress of Semna records an anonymous appealing to Queen Katimala for help against rebel chiefs \parencites{darnellInscriptionQueenKatimala2006}[p.~296]{törökTwoWorldsFrontier2009}. Darnell assigned inscription to approximately 1069---720 BCE. He argues that its content points to a fractured Nubia, makerd by hostilities between "the followers of Amun and the marauders from the desert." Furthermore, he suggests that the inscription's proximity to gold deposits implies that control over this resource was a key driver of these conflicts \parencite[p.~60]{darnellInscriptionQueenKatimala2006}. However, other scholars have criticized the dating, proposing a date closer to the 25th Dynasty (ca. 744---656 BCE) \parencite[p.~298]{törökTwoWorldsFrontier2009}. At the current state of knowledge, Queen Katimala's inscription may or may not be a clue as to the political organization of Nubia during the post-New Kingdom period and the eventual rise of the Napatan kings.


Nonetheless, by the 8th century BCE, a unified kingdom had been established in Upper Nubia which would eventually expand north and reach Thebes. During the 25th Dynasty, evidence written in the Ancient Egyptian language, along with Manetho's list of kings, provided scholars with records of the names of the Kushite rulers \parencites{payraudeauRetourSuccessionShabaqoShabataqo2014}{pope25thDynasty2019}.\footnote{Scholars have proposed different absolute dates for the 25th Dynasty rulers \parencites{depuydtDatePiyesEgyptian1993}{kahnInscriptionSargonII2001}{hornungKushiteRulersPre25th2006}{payraudeauRetourSuccessionShabaqoShabataqo2014}. For this thesis, I will employ Payraudeau's suggested chronology.}


The first documented Kushite ruler is Alara, whose name only appears in inscriptions of his successors \parencite[p.~110]{emberlingKushDynastyNapata2023}. Alara was succeeded by Kashta\footnote{Although the name Kashta has been interpreted to mean "the Kushite", \textcite{vinogradovDidNameKashta2003} demonstrates that the words "Kush" and "Kashta" follow distinct spelling logic.}, who installed his daughter, Amenirdis I, as God's Wife of Amun at Thebes, establishing his control over Upper Egypt \parencites[p.~319]{törökTwoWorldsFrontier2009}[p.~110]{emberlingKushDynastyNapata2023}.\footnote{On God's Wives of Amun during the 25th Dynasty, see \cite{ayadGenderRitualManipulation2016}, \cite{beckerFemaleInfluenceAside2016}, \cite{lohwasserNubianessGodsWives2016}.} 
His successor, Piankhy (ca. 744---714 BCE), adopted the traditional Egyptian five-part royal titles \parencites[p.~153]{törökKingdomKushHandbook1997}[p.~324]{törökTwoWorldsFrontier2009}.\footnote{The name Piye is also commonly used in the literature; however, \textcite{rillyNouvelleInterprétationNom2001} argues in favour of using Piankhy.} 
In his "Victory Stela" (currently at the Egyptian Museum in Cairo, JE 48862), likely erected in Year 3 of his reign, Piankhy declared himself the legitimate ruler of Egypt and intended to keep Egypt's political map fragmented as long as the other rulers recognized his supremacy and payed him tribute \parencite[p.~155]{törökKingdomKushHandbook1997}. The following kings, Shabataqo and Shabaqo,\footnote{It was previously believed that Shabaqo was followed by Shabataqo. Recent scholarship suggests reversing the succession order, supported by Kushite and Egyptian textual evidence \parencites[p.~4-9]{payraudeauRetourSuccessionShabaqoShabataqo2014}{jurmanOrderKushiteKings2017}, as well as analysis of horse burials at El-Kurru \parencite{näserKingsHorsesTwo2020}.} relocated the royal court to Memphis. Shabaqo further consolidated the double kingdom by controlling the local rulers of the Delta region \parencites[p.~166-167]{törökKingdomKushHandbook1997}[p.~141]{rillyHistoireSoudanOrigines2017}.\footnote{A special case is Tanis, where no local rulers were documented and was ruled directly by Shabaqo \parencite{mojeHerrschaftsräumeUndHerrschaftswissen2013}.}


Shabaqo was suceeded by Taharqa, who engaged in building and enlarging temples in Egypt and Kush (ca. 690---664 BCE).\footnote{For a more detailed study on Taharqa's double kingdom, see \cite{dalliborTaharqoPharaoAus2005}, \cite{popeDoubleKingdomTaharqo2014}.} However, three years after successfully driving out Assyrian troops in Palestine, Memphis was invaded and sacked, forcing Taharqa to flee to Thebes. The Assyrians installed new rulers---both Assyrian and Egyptian---in Lower Egypt, effectively controlling the whole region. Although Taharqa briefly regained control of Memphis, his success was short-lived. By approximately 664 BCE, the Assyrians reasserted dominance over Lower Egypt and decisively defeated the Kushite-Egyptian forces in Upper Egypt. Following this defeat, Taharqa permanently retreated to Napata \parencite[p.~145-154]{rillyHistoireSoudanOrigines2017}.


His successor, Tanwetamani, claims to have briefly regained control of Lower and Upper Egypt, before being driven away to Kush when the Assyrians conquered Upper Egypt again, as well as looting and burning Thebes. The double kingdom was effectively over, even though commercial and trade relations were maintained between Egypt and Kush \parencite[p.~356-360]{törökTwoWorldsFrontier2009}. Tanwetamani was buried at El-Kurru and  not in the new established cemetery at Nuri.\footnote{Only one tomb is thought to have been built at El-Kurru again, pyramid Ku. 1. Since the tomb had no inscribed material, it is not known who was the supposed owner. Recent excavations agree that it was build in the late Napatan or early Meroitic periods and was left unfinished. The excavators tentatively interpret that this pyramid remained unused due to political conflicts among rulers in Napata. These conflicts could have led to the transition from Nuri as the royal burial site to Meroe at the end of the 4th century BCE \parencites[p.~62]{emberlingExcavationPyramidKu2015}{yellinPyramidChapelDecorations2015}.}


The following four kings---Atlanersa, Senkamanisken, Anlamani, and Aspelta---, collectively reigned for approximately eighty years. During their reign, they built and expanded temples in Nubia, commissioned statues of themselves and earlier rulers, such as Taharqo and Tanwetamani, and issued public inscriptions and seals \parencites[p.~157-163]{rillyHistoireSoudanOrigines2017}[p.~133-137]{emberlingKushDynastyNapata2023}. Around 593 BCE, Psamtek II, king of Egypt, invaded Kush. Most scholars believe this invasion occurred during Aspelta's reign, as the sequence of destroyed Kushite royal statues ends with him. The burning of the “Treasury” at Sanam is also considered a consequence of this raid \parencites[p.~361-362]{törökTwoWorldsFrontier2009}{andersonWhatAreThese2009}[p.~138-139]{emberlingKushDynastyNapata2023}.


Apart from their names and tombs at Nuri, very little is known about the approximately seventeen succeeding kings \parencites[p.~435]{williamsNapatanNeoKushiteState2020}[p.~139]{emberlingKushDynastyNapata2023}.\footnote{For the complete list see \cite[p.~496-497]{hornungKushiteRulersPre25th2006}.} They are followed by a group of five rulers known as “Neo-Ramesside kings” due to their titles and inscriptions bearing similarities to the Egyptian Ramesses. The dating of these rulers has been a subject of debate. Some scholars have placed them as the last kings of the Napatan period due to the paleography and contents of their titles \parencite[p.~293]{törökTwoWorldsFrontier2009}. Whereas others have argued that they should be placed before the 25th Dynasty, since the Libyan kings of the 23rd Dynasty also emulated Ramesside titles \parencites[p.~216-217]{morkotEmptyYearsNubian1991}[p.~145-150]{morkotBlackPharaohsEgypts2000}[p.~63-64]{kendallOriginNapatanState1999}. Recently, Lohwasser published a comprehensive reassessment of the textual evidence concerning the Neo-Ramesside kings, tentatively concluding that only two of these kings could possibly be attributed to the end of the Napatan period \parencite[p.~132-134]{lohwasserKushiteKingsTwentyfifth2022}. Therefore, the dating of these kings remains uncertain due to the limited and inconclusive nature of the available evidence.


Regardless of who were the last kings of the period, the primary event signaling the end of the Napatan dynasty was the relocation of the royal cemetery from Nuri to Meroe with the burial of Arkamani I around 270 BCE, effectively shifting the political and economic center from Napata to Meroe. The transition to the Meroitic period---which lasted until ca. 350 AD--- also involved cultural developments that were gradually emerging. Among those are the cult of local deities, specially the lion god Apedemak, the introduction of ruling queens known as \textit{kandake}, and the use of the Meroitic language \parencites[p.~144]{elhassanRoleEarlyNapatan2023}[p.~2]{kuckertzMeroeEgypt2021}.


\section{Tracing the Napatan dynasty: chronological and emergence issues}


The origins of the Napatan rulers has been the subject of debate for many years.\footnote{For a comprehensive overview of the literature on the many theories regarding the Napatan dynasty origins and initial chronology, see \cite{morkotPriestlyOriginNapatan2003}.} In the centre of the discussions is the cemetery of El-Kurru, considered as the burial site of their the ancestors \parencite{törökKingdomKushHandbook1997, törökTwoWorldsFrontier2009, morkotKingshipKinshipEmpire1999}. The importance of El-Kurru is twofold, involving the emergence problem and the overall chronology of the period, which affects the ruling dates of the kings. As Török notes, "the different el-Kurru chronologies describe and at the same time create different images of political and religious processes" \parencite[p.~304]{törökTwoWorldsFrontier2009}. Scholars have suggested different models for the emergence of the Napatan dynasty and chronology of el-Kurru---known as the short and long chronology.


The history of the cemetery was first established by its excavator Reisner and published by Dunham in 1950. Dunham proposed a chronology based on Reisner's, with some of his own modifications \parencite[p.~2]{dunhamElKurru1950}. Reisner identified sixteen tombs and divided them into six different types. Starting at around 860---840 BCE, the tombs would, therefore, be of six different generations of rulers and their wives \parencite{reisnerDiscoveryTombsEgyptian1919, reisnerNoteHarvardBostonExcavations1920}. Known as the short chronology, this division was adopted by Kendall, who---drawing on new evidence---assigned the earlier tombs to around 885---835 BCE and revised the burial sequence \parencite[p.~4]{kendallOriginNapatanState1999}. Regarding the origin of the Napatan rulers, Kendall argues that a new influx of Egyptian settlers was a key influence. He bases this conclusion on the correlation between Egyptian influences at the El-Kurru burials and the rebellion of Theban priests against the new High Priest of Amun, around 839 BCE. Kendall suggests that these priests fled southward and subsequently entered under the service of the Nubian kings \parencite[p.~5,57]{kendallOriginNapatanState1999}.  Therefore, according to the author, the El-Kurru ancestral rulers were a non-Egyptianized Nubian group that, through the external influence of native Egyptians, were rapidly "Egyptianized".


Kendall's model was criticized by Török, who proposed a longer chronology \parencite[p.~298-309]{törökEmergenceKingdomKush1995, törökOriginNapatanState1999, törökTwoWorldsFrontier2009}. He argues for an earlier date for the first burials, around 995---975 BCE, and asserts that all central sixteen burials belonged to ruling males \parencite[p.~301]{törökTwoWorldsFrontier2009}. Furthermore, Török suggests that the Nubian elite, previously integrated into the New Kingdom Egyptian administration, maintained their power after the withdrawal, taking advantage of the direct access to gold mining sites \parencite[p.~111-112]{törökKingdomKushHandbook1997}. The author characterizes the Napatan state as a transformation from "chiefdom" to kingdom, with its Egyptianized features resulting from the fusion of indigenous and Egyptian concepts \parencites{törökKingdomKushHandbook1997, törökTwoWorldsFrontier2009}.\footnote{Török does not explicitly define "chiefdom", but appears to distinguish it from "kingdom" based on territorial extent, using the latter to denote Napatan rule over the whole Middle Nile region \parencite[p.~193]{törökKingdomKushHandbook1997}} Other scholars have also supported this theory on the origins of the Napatan dynasty \parencites{morkotFoundationsKushiteState1995, morkotEgyptNubiaEgyptian2001}{yellinEgyptianReligionIts1995}{smithNubiaEgyptInteraction1998}{binder10th9thCenturyBC2011}. Table \ref{tab:el-kurru-tombs} presents the short and long chronologies of the ancestral El-Kurru burials alongside their proposed occupants.


{\small
	\begin{xltabular}{\textwidth}{X X X X X X}
		\caption{El-Kurru ancestral tombs with their respective dating and occupant according to the short and long chronology}
		\label{tab:el-kurru-tombs} \\
		\toprule
		\multicolumn{3}{c}{\textbf{Kendall's short chronology (1999)}} & \multicolumn{3}{c}{\textbf{Török's long chronology (2009)}} \\
		\cmidrule(r){1-3} \cmidrule(l){4-6}
		\textbf{Tomb} & \textbf{Dating} & \textbf{Owner} & \textbf{Tomb} & \textbf{Dating} & \textbf{Owner} \\
		\midrule
		\endfirsthead
		\multicolumn{6}{c}{\textit{Continued from previous page}} \\
		\toprule
		\multicolumn{3}{c}{\textbf{Kendall's short chronology (1999)}} & \multicolumn{3}{c}{\textbf{Török's long chronology (2009)}} \\
		\cmidrule(r){1-3} \cmidrule(l){4-6}
		\textbf{Tomb} & \textbf{Dating} & \textbf{Owner} & \textbf{Tomb} & \textbf{Dating} & \textbf{Owner} \\
		\midrule
		\endhead
		\bottomrule
		\endlastfoot
		Ku. Tum. 1 & ca. 885-835 BCE & Lord A & Ku. Tum. 1 & ca. 995-975 BCE & Prince \\
		\cmidrule(r){1-3} \cmidrule(l){4-6}
		Ku. Tum. 5 & ca. 885-835 BCE & Wife of Lord A & Ku. Tum. 4 & ca. 975-955 BCE & Prince \\
		\cmidrule(r){1-3} \cmidrule(l){4-6}
		Ku. Tum. 4 & ca. 865-825 BCE & Wife of Lord B & Ku. Tum. 5 & ca. 955-935 BCE & Prince \\
		\cmidrule(r){1-3} \cmidrule(l){4-6}
		Ku. Tum. 2 & ca. 865-825 BCE & Lord B & Ku. Tum. 2 & ca. 935-915 BCE & Prince \\
		\cmidrule(r){1-3} \cmidrule(l){4-6}
		Ku. Tum. 6 & ca. 845-815 BCE & Lord C & Ku. Tum. 6 & ca. 915-895 BCE & Prince \\
		\cmidrule(r){1-3} \cmidrule(l){4-6}
		Ku. 19 & ca. 845-815 BCE & Wife of Lord C & Ku. 19 & ca. 895-875 BCE & Prince \\
		\cmidrule(r){1-3} \cmidrule(l){4-6}
		Ku. 14 & ca. 825-805 BCE & Lord D & Ku. 14 & ca. 875-855 BCE & Prince \\
		\cmidrule(r){1-3} \cmidrule(l){4-6}
		Ku. 13 & ca. 825-805 BCE & Wife of Lord D & Ku. 13 & ca. 855-835 BCE & Prince (?) \\
		\cmidrule(r){1-3} \cmidrule(l){4-6}
		Ku. 11 & ca. 805-795 BCE & Lord E & Ku. 11 & ca. 835-815 BCE & Prince \\
		\cmidrule(r){1-3} \cmidrule(l){4-6}
		Ku. 10 & ca. 805-795 BCE & Wife of Lord E & Ku. 10 & ca. 815-795 BCE? & Unknown \\
		\cmidrule(r){1-3} \cmidrule(l){4-6}
		Ku. 9 & ca. 785 BCE & Alara & Ku. 9 & ca. 795-775 BCE & Alara \\
		\cmidrule(r){1-3} \cmidrule(l){4-6}
		Ku. 23 & ca. 785 BCE & Queen Kasaqa & Ku. 8 & ca. 775-755 BCE & Kashta \\
		\cmidrule(r){1-3} \cmidrule(l){4-6}
		Ku. 21 & ca. 760 BCE & Minor queen (?) & Ku. 17 & ca. 755-721 BCE & Piankhy \\
		\cmidrule(r){1-3} \cmidrule(l){4-6}
		Ku. 8 & ca. 760-747 BCE & Kashta & & & \\
		\cmidrule(r){1-3} \cmidrule(l){4-6}
		Ku. 7 & ca. 747-736 BCE & Queen Pebatma & & & \\
		\cmidrule(r){1-3} \cmidrule(l){4-6}
		Ku. 20 & after ca. 747 BCE & Minor queen (?) & & & \\
		\cmidrule(r){1-3} \cmidrule(l){4-6}
		Ku. 17 & ca. 716 BCE & Piankhy & & & \\
	\end{xltabular}
}

Recent excavations at the Tombos cemetery in Upper Nubia support this model, as there is little evidence of a new wave of Egyptian migration following the New Kingdom \parencite{buzonEntanglementFormationAncient2016}. In addition, a horse burial from the cemetery---dated to approximately 1005---851 BCE---, is interpreted as a symbol of emerging statehood, since it falls between the scattered horse burials of the New Kingdom and the more ritualized horse burials at El-Kurru, associated with the kings of the 25th Dynasty. As a result, this feature of Napatan culture seems to have originated in smaller communities and was later incorporated into the state \parencite{schraderColonialindigeneInteractionAncient2018}.

At present, and with the available evidence, the exact dating and the origins of the first Napatan rulers cannot be absolutely determined. However, since Nubia was incorporated into the political and economic structure of the New Kingdom Egyptian administration through the authority of local elites, the Napatan kingdom could be seen as a development of the increasing power of these elites. Morkot named this phenomenon "post-imperial" or "post-colonial", as the roles between the two societies were reversed \parencite[p.~914-915]{morkotConqueredConquerorOrganization2013}. In this sense, this hypothesis not only takes into account the authority of the indigenous elites in the aftermath of the New Kingdom occupation but also highlights Nubian agency in controlling their own sociopolitical landscape.


\section{Cultural interactions between groups within Nubia and Egypt}


Historical exchanges between Nubia and Egypt have significantly influenced the development of both civilizations, with evidence of contact attested from at least the 4th millennium BCE onward. These interactions were facilitated by extensive trade networks, through which goods and people moved between regions, while periods of conflict often resulted in shifts in power dynamics. During periods of unification, the Egyptian state was interested in the Nubian region to exploit its natural and human resources, and territorial expansion \parencites[p.~94-95]{seidlmayerNubierImÄgytischen2002}[p.~570]{raueNubiansEgypt3rd2019}. Furthermore, throughout periods of instability and fragmented power---e.g. First and Second Intermediate Periods---, local rulers in Egypt turned towards Lower Nubia in search of trade agreements as well as soldiers to serve their own interests (although there were also periods of conflicts with Kushite powers) \parencites{daviesKushEgyptNew2003}[p.~141]{näserStructuresRealitiesEgyptianNubian2013}{polzTerritorialClaimPolitical2018}{morenogarcíaShapingNewIdentities2024}.


By examining these complex interactions, we can better understand how these societies mutually shaped each other and contributed to the cultural composition of the region. In this section, a brief overview of Nubian-Egyptian interactions through material culture is first provided, followed by a specific focus on the New Kingdom and Napatan period, highlighting two distinct modes of interaction between these societies.


The earliest Nubian society known to have had contact with Egypt, the A-Group, occupied Lower Nubia and dated to approximately 3700---2800 BCE. Burials included mud brick covering and Egyptian imports, such as vessels, beads, amulets, seals and combs \parencites{gattoNubianAGroupReassessment2006}[p.~35-36]{törökTwoWorldsFrontier2009}[p.~83]{williamsRelationsEgyptNubia2011}[p.~184-199]{royPoliticsTradeEgypt2011}[p.~294]{raueCulturalDiversityNubia2019}[p.~134, 136]{gattoAGroup4thMillennium2020}. Additionally, a seal found south of the First Cataract was decorated both with Egyptian royal symbols and Nubian elements \parencite[p.~134, 136]{gattoAGroup4thMillennium2020}. At the same time, A-Group objects are also attested at Egyptian Predynastic sites \parencites{gattoNubiansEgyptSurvey2005}[p.~63-64]{gattoNubianAGroupReassessment2006}{gattoEgyptNubia5th4th2009}[p.~95]{gattoCulturalEntanglementDawn2014}. Consequently, throughout Upper Egypt, both cultures coexisted and exhibited a blend of Egyptian and Nubian customs \parencites[p.~129]{gattoEgyptNubia5th4th2009}[p.~290-291]{meurerNubiansEgyptEarly2020}. At Elephantine and Buhen, the contact between Nubians and Egyptians has been called "Dynastic A-Group" \parencite[p.~4]{raueWhoWasWho2008}. It has been proposed that long-distance trade between Upper Egypt and the A-Group, with the introduction of prestige goods, played a role in the emergence of the unified Egyptian state under the 1st Dynasty (ca. 2900 BCE) \parencite[p.~49]{törökTwoWorldsFrontier2009}. By the end of the 4th millennium BCE, an unified Egypt state launched military campaigns into Lower Nubia, effectively bringing about the disappearance of the A-group culture \parencites{hafsaasEncounteringExpandingPredatory2026}.


Following the A-Group, the C-Group occupied Lower Nubia from around 2500 to 1500 BCE. Evidence of long-distance trade with Egypt is attested by grave goods and graffiti left by caravan leaders under king Pepy I \parencites[p.~166]{hafsaasCGroupPeopleLower2020}. During the First Intermediate Period, as Egyptian groups migrated to Lower Nubia, the influx of imports increased, including amulets, metal objects, scarabs, and pottery jars filled with grains \parencites[p.~303]{raueCulturalDiversityNubia2019}[p.~171]{hafsaasCGroupPeopleLower2020}. In Egypt, tomb inscriptions and other written records provide insights into the political landscape of Nubia, naming its rulers and detailing expeditions to quarries. C-Group people were part of Egyptian society, being employed as mercenaries, in royal court and domestic contexts \parencites[p.~96-99]{seidlmayerNubierImÄgytischen2002}[p.~569]{raueNubiansEgypt3rd2019}. Settlements have been discovered near Elephantine, as well as a cemetery at Hierakonpolis, marking the northernmost extent of known C-Group presence in Egypt. The burials showcase a blend of Nubian and Egyptian traditions \parencites{friedmanNubianCemeteryHierakonpolis2007}[p.~292]{meurerNubiansEgyptEarly2020}.


During the Middle Kingdom, Egypt occupied Lower Nubia, and although the C-Group initially resisted, they were eventually subjugated and collaborated with their new rulers \parencites[p.~172]{hafsaasCGroupPeopleLower2020}. This conquest allowed Egypt to acquire resources directly, reducing the need for trade and imposing stricter control between the various Nubian communities. However, the Egyptian presence in the region also benefited the local trade \parencite[p.~138]{näserStructuresRealitiesEgyptianNubian2013}. In royal contexts, three wives of the 12th Dynasty king Mentuhotep II have been considered by some scholars as Nubian due to their representation as dark-skinned on funerary reliefs as well as their tattoos \parencites[p.~67]{ashbyDancingHathorNubian2018}[p.~572-573]{raueNubiansEgypt3rd2019}[p.~71-73]{renautTattooedWomenNubia2020}


While the Egyptian centralized government collapsed during the Second Intermediate Period, other cultures emerged in Nubia, such as the Kingdom of Kerma in Upper Nubia and the Pan-Grave culture in Middle Nubia. Amidst hostilities between Thebes and Kerma during the end of the Second Intermediate Period (ca. 1555---1550 BCE), the C-Group aligned with the northern power. As a consequence, it has been suggested that the local elite adopted more Egyptian cultural markers, and even held positions of power under the Egyptian viceroy \parencites{hafsaas-tsakosKushEgyptCGroup2010}{hafsaasCGroupPeopleLower2020}.


Pan-Grave communities were also in contact with Egyptians. They were a nomadic group believed to have inhabited the Eastern Desert during the 2nd millennium BCE and their presence becomes visible when they began interacting with Egypt and other Nubian groups \parencites[p.~81, 89]{näserNomadsNileArchaeology2012}[p.~175]{souzaNotMarginalMarginalised2022}. Evidence of interactions is attested through the adoption of Egyptian burial customs, such as the extended burial position, and the presence of their pottery in Egyptian settlements  \parencites[p.~294-296]{meurerNubiansEgyptEarly2020}[p.~4]{desouzaMeltingPotsEntanglement2020}[p.~125-126]{desouzaBlurringBoundariesMore2024}. However, their role in Egyptian society is still a matter of debate
\parencite[p.~229]{liszkaPanGraveMedjayIntersection2020}.


Between 2500 and 1500 BCE, the Kerma kingdom founded a city with the same name in Upper Nubia. In Egypt, elite burials contained numerous Kerma objects. Notably, the treasure of New Kingdom Queen Ahhotep featured three fly amulets crafted from Nubian gold, which became highly popular during the 18th Dynasty \parencite[p.~188]{williamsKushWiderWorld2020}. Small Kerman cemeteries are known as far north as Saqqara, but no settlements have been found. Employment of Nubian mercenaries by the Egyptian state are suggested by Nubian material culture found in a military settlement at Deir el-Ballas \parencite[p.298-300]{meurerNubiansEgyptEarly2020}.


Conversely, Kerma tumulus tombs exhibit a blend of local and Egyptian elements, including seal impressions from the Middle Kingdom and Second Intermediate Period \parencites[p.~190-191]{williamsKushWiderWorld2020}[p.~205-206]{bonnetCitiesKermaPnubsDokki2020}. Over time, Egyptian imports become increasing common in Kerma burials and are especially high in royal tombs \parencite[p.~216-219]{bonnetEasternCemeteryKerma2020}. Elite graves, on the other hands, showcase traditional Kerma features, such as bed burials and animal or human sacrifices, as well as local made objects incorporating Egyptian styles \parencite[p.~118]{emberlingEarlyKushKingdom2022}. Additionally, strontium analysis of human remains from the Abu Fatima cemetery and the city of Kerma indicate a mix of local and non-local individuals during this period, with a higher proportion of non-locals at Abu Fatima (24\%) compared to Kerma (13\%) \parencite[p.~377]{schraderIntraregional87Sr86Sr2019}.


It is during this period---around 2000---1750 BCE--- that the first written evidence of the name “Kush” appears in Egypt. As Kerma expanded to the north, this led to conflicts between both societies, leading Egypt to build several fortresses in Lower Nubia during the Middle Kingdom (ca. 1985---1773 BCE) to control trade, people and the Nubian expansion \parencites[p.~130-135]{emberlingPastoralStatesComparative2014}[p.~230]{honeggerHolocenePrehistoryUpper2019}[p.~219-220]{bonnetEasternCemeteryKerma2020}[p.271]{bestockEgyptianFortressesColonization2020}[p.~130]{emberlingEarlyKushKingdom2022}.


By the mid-second millennium, a now unified and imperialistic Egypt asserted control first over Lower Nubia and eventually to Upper Nubia, destroying the city of Kerma \parencite[p.~78]{edwardsNubianArchaeologySudan2004}. Egyptian control was extended as far as the 4th Cataract, giving it access to Nubia's gold mines and trade routes. By the reign of Thutmose I (1506---1493 BCE), local rulers were integrated into the Egyptian administration \parencites[p.~76]{morkotBlackPharaohsEgypts2000}[p.~99]{törökKingdomKushHandbook1997}. According to \textcite[p.~103-104]{morrisAncientEgyptianImperialism2018}, the New Kingdom colonialism relied on different strategies: it aimed to transform soldiers into settlers, local elites into "Egyptianized" rulers, and the indigenous religious sphere into the cult of Amun, at the same time as reshaping local governance and economy to reflect the Egyptian model. This meant that the Nubian trade network was incorporated into the Egyptian redistributive system \parencites[p.~157]{törökKingdomKushHandbook1997, törökTwoWorldsFrontier2009}{morkotConqueredConquerorOrganization2013}. Egypt organized the region into Wawat (Lower Nubia) and Kush (Upper Nubia), each governed by viceroys, with temple towns providing a new venue of interactions between both cultures, with the production, consumption and importing of goods, especially from Egypt \parencites{kempFortifiedTownsNubia1972}{roseSesebiCeramicsChronology2017}[p.~343]{spencerBuildingNewGround2017}{viethNewKingdomTemple2019}{budkaMaterialRemainsNew2020}. According to \textcite{smithWretchedKushEthnic2003}, these towns were inhabited by indigenous peoples, Egyptians, and those of mixed heritage through intermarriage. 


Important temples dedicated to Amun were built in cities such as Napata, Kerma, and Kawa. These temples were home to various manifestations of Amun, including a ram-headed form, which the Nubians regarded as their local version of the deity \parencite{gaboldeAmunCultIts2020}. It has been proposed that this form of Amun represents a fusion of the Egyptian god with a local ram cult from the Kerma culture \parencites[p.~77]{bonnetKermaRoyaumeNubie1990}{valbelleAmonRêKerma2003}[p.~227]{törökTwoWorldsFrontier2009}. Gabolde and Török hold differing views on the extent of the Amun cult's integration into Nubian society. Török argues that the temple iconography---accessible to the general populatios---suggests that lower social groups were also engaged with the cult \parencite[p.~210]{törökTwoWorldsFrontier2009}, while Gabolde contends that the cult's influence was primarily confined to the elite \parencite[p.~354]{gaboldeAmunCultIts2020}.


Archaeological data from Nubian towns provide insights into the cultural entanglement occurring there. For instance, the presence of local architecture, the use of Nubian handmade pottery alongside Egyptian wheelmade vessels, and Egyptian-style female figurines with distinct Nubian features (decoration, pronounced hips) \parencites{spencerNubianArchitectureEgyptian2010, spencerIntroductionHistoryHistoriography2017, spencerMaintainingRamessideEmpire2024}[p.~183]{budkaLifeNewKingdom2012}[p.~471]{roseSesebiCeramicsChronology2017}{stevensFemaleFigurinesFolk2017}[p.~225]{budkaMaterialRemainsNew2020}
At the settlement of Askut, the presence of Nubian cooking pots within domestic contexts indicates that Egyptian migrants married Nubian women, and that these women maintained certain local traditions \parencites{smithWretchedKushEthnic2003, smithEthnicityCulture2007}[p.~540]{vanpeltRevisingEgyptoNubianRelations2013}.


Cemeteries contribute to understand the cultural landscape. The burial of a local individual from Amara West blended the Nubian tumulus with the Egyptian underground chamber and extended body position. Additionally, grave goods were of Nubian origin, differing from the contemporary Egyptian practices by focusing on daily life items rather than traditional funerary objects \parencites[p.~594-607]{binderNewKingdomTombs2017}[p.~44-45]{spencerMaintainingRamessideEmpire2024}.
Similarly, at Sai Island, an Egyptian-style pyramid and shaft tomb belonging to Khnummose, the Overseer of Goldsmiths, contained Egyptian grave goods even though isotopic analysis indicated that he originated from Nubia \parencite{retzmannNewKingdomPopulation2019}. Research from the Tombos cemetery have demonstrated the different levels of entanglement: while locals were buried following customs of both cultures, non-locals were employed only Egyptian customs. This evidences the flexibility and agency in Nubian funerary choices rather than adhering to a single “Egyptianized” tradition \parencites[p.~6]{buzonStrontiumIsotope87Sr2013}[p.~291-292]{buzonEntanglementFormationAncient2016}[p.~207]{smithFortifiedSettlementTombos2018}{schraderIntraregional87Sr86Sr2019} [p.~102]{smithCrossCulturalInteractionAncient2018} [p.~12-16]{buzonExploringIntersectionalIdentities2024}.


Local elite tombs also offer additional perspectives beyond material culture and human remains. An example is the tomb of Hekanefer, who served the Egyptian New Kingdom administration as “Prince of Aniba”. In his tomb, he is depicted in an Egyptianized style \parencite{simpsonHekaneferDynasticMaterial1963}. Interestingly, Hekanefer also appears in the tomb of Huy in Thebes, the Viceroy of Kush, presenting tribute to King Tutankhamun. In this depiction, he wears Egyptian dress while retaining Nubian cultural markers—ostrich feathers on his head, penannular bracelets, earrings, and animal skin. The diverging depictions serve two purposes. On one hand, an Egyptianized depiction in a Nubian elite tomb signified higher status, prestige, and integration into the administration. On the other hand, in a foreign context, Nubian elites asserted their identity through distinct cultural markers, maintaining a balance between assimilation and self-representation \parencites{törökTwoWorldsFrontier2009}{vanpeltRevisingEgyptoNubianRelations2013}{smithHekaneferLowerNubian2015}.


Lemos' analysis of grave goods in elite and non-elite cemeteries reveals significant differences in how identity and status were negotiated in New Kingdom Nubia. Elite burials, such as those of Khnummose at Sai Island and Weser in Aniba, demonstrate access to restricted Egyptian objects, such as shabtis and two heart scarabs, and how locally available jars were repurposed as "canopic jars," demonstrating how imported and local objects were adapted and transformed within the burial context \parencites[p.~7]{lemosAlternativesColonizationMarginal2021}[p.~76-77]{lemosCulturalEntanglementsExperiencing2024}. Furthermore, by examining patterns—or the absence thereof—in the distribution of grave goods across burials at Aniba, Sai, Soleb, and Fadrus, Lemos revealed  internal diversity within the burial assemblages. This diversity suggests that cultural entanglements were not uniform but instead shaped by local contexts. Such variations may reflect not only regional preferences but also constraints related to the consumption of specific objects by particular social groups \parencite{lemosForeignObjectsLocal2020}.


Following the Egyptian withdrawal, archaeological evidence indicates that temple towns and cemeteries remained inhabited, where interactions between Egyptians and Nubians residents persisted even after the colonial administration disappeared. The cemetery of Tombos also remained in use during the Napatan period. Egyptian-style burials, featuring pyramids and underground chambers, continued alongside Nubian-style tumuli. Most burials followed the Egyptian extended position, and grave goods included Egyptian-style objects, such as an inscribed heart scarab found near a pyramid and a scarab bearing the name of Piankhy \parencite[p.~293-294]{buzonEntanglementFormationAncient2016}. Traditional Nubian objects include Red Sea shell beads, ivory bracelets and quartz crystal pendants that resemble those dated to the Kerma Period. Particularly interesting objects include a headrest similar to those found in Kerma burials---which were already adaptations of typical Egyptian headrests---and a heart scarab that showcases unique artistic practices, referencing Egyptian styles but not directly imitating them. In addition, the pottery showcases a fusion of Nubian and Egyptian influences, combining decorative motifs and manufacturing techniques from both cultures \parencites[p.~6-9]{smithColonialEntanglementsEgyptianization2014}[p.~637-642]{buzonTumuliTombosInnovation2023}.


The tumuli typically held one individual, a practice linked to burial traditions at Kerma. In these burials, Egyptian-style coffins were often paired with Nubian burial beds. Isotopic analysis of the human remains revealed that none of the Napatan samples had Egyptian strontium values, indicating that Egyptian immigration had ceased by this time. However, morphological and isotopic evidence suggests that Nubians from other regions were integrated into the Tombos community. Researchers propose that this population may represent descendants of the mixed Nubian-Egyptian community from the New Kingdom \parencites{schraderIlluminatingNubianDark2014}[p.~294-295]{buzonEntanglementFormationAncient2016}{buzonTumuliTombosInnovation2023}{buzonExploringIntersectionalIdentities2024}{smithCrossCulturalInteractionAncient2018}.


Similar to Tombos, the burials from at Sedeinga feature both tumuli and pyramids \parencites{berger-elnaggarContributionSedeingaLhistoire2008}{rillyCloserAncestorsExcavations2018}. Recent excavations have revealed a particularly interesting burial. The undisturbed grave II TT 262 belonged to a very young individual, less than two years old, who was interred in a lateral niche grave topped by a tumulus superstructure. Interestingly, this grave was related to two pyramids, as it was constructed behind them. This placement highlights the blend of local tumulus traditions with the architectural style of Egyptian pyramids. The grave goods included an Egyptian-style amulet depicting the god Shu and copper-alloy anklets. Ceramic assemblages from other graves comprise both Egyptian imports and local productions \parencites{rillySedeinga2012Season2013, rillyCloserAncestorsExcavations2018, rillyCollectiveGravesBastatues2020}. Similar evidence from pottery is also attested at the of Kirbekan in the region of the 4th Cataract \parencites{budkaDocumentationExcavationDome2007}. The transition cemeteries of Amara West---Cemeteries C and D---feature tombs spanning from the New Kingdom to the early Napatan period. This site, like the other cemeteries in the region, exhibits a similar blend of Nubian and Egyptian burial traditions. Egyptian chamber tombs, for instance, incorporated both Egyptian elements, such as headrests and coffins, and Nubian cultural features, including burial beds and flexed body positions. Similarly, Nubian tumuli tombs often contained coffins and extended burials. Both types of burials contained Egyptian-style amulets \parencites{spencerAmaraWestII2002}{binder10th9thCenturyBC2011}.


During the Kushite 25th Dynasty in Egypt, depictions of Kushites exhibited Egyptian, Nubian, or a blend of both features. In some cases, it is only through their father's name that we recognize them as Nubian, suggesting that these individuals had the flexibility to adapt or change their identity as they wished \parencites{budkaKushiteTombGroups2010, budkaNubians1stMillennium2019, budkaNubiansEgypt25th2020}. At Abydos, five monuments attest to the presence of Nubians, including a stela depicting the Kushite prince Ptahmaakheru, who likely belonged to the royal family. Other monuments feature another member of the royal family and an official or member of the royal household \parencite{budkaKuschitenAbydosEinige2012}{leahyKushitesAbydosRoyal2014}. These can be identified by their non-Egyptian name, costume and physical depiction \parencite[p.~704-705]{budkaNubians1stMillennium2019}. 


The Theban region included a cluster of Kushite tombs, as well as intrusive Kushite burials in Tomb TT 99 of Sennefer. This cluster included the coffin of Niw, daughter of Padiamun, Priest of Amun in the Nubian city of Kawa. Her depiction reveals a blend of Egyptian and Nubian features, such as an Egyptian garment paired with a short Kushite hairstyle \parencite[p.~504-505]{budkaKushiteTombGroups2010}. Another tomb, featuring a mud brick superstructure, contained the remains of several individuals of Kushite ancestry. Although much of the burial assemblage reflects Egyptian characteristics, certain distinct Kushite elements were preserved. These include personal names, depictions, Nubian objects, and Egyptian grave goods with unique Nubian attributes. The apparent absence of Kushite artifacts in Egypt could be due to either a lack of resources to personalize Egyptian objects or a lack of interest in doing so, which contributed to a uniformity in the archaeological record \parencite[p.~494]{howleyKushitesEgypt6642020}. While archaeological evidence for Nubians—especially from lower social groups—is scarce, it is likely that they settled in Egypt during the 25th Dynasty. Additionally, some may have married Egyptians, leading to the formation of intercultural families similar to those found in Nubian towns \parencite[p.~477]{budkaNubiansEgypt25th2020}.


In Nubia, royal burials also contain blended material culture. A re-examination of shabtis from three queen tombs at El-Kurru (Ku. 62, 71 and 72) uncovered a particular type featuring a basket on its head, which only parallel comes from an elite New Kingdom tomb in Thebes (TT 99) \parencite{mussoKushiteShabtisBasket2011}. Imported marl pottery jars from the same site indicate trade connections with Upper Egypt, which apparently ceased during the later half of the 7th century BCE \parencite{heidornBostonMuseumFine2018}. Additionally, pot-breaking ceremonies observed in tomb enclosures at El-Kurru mirror those performed by Egyptian priests at Abydos \parencite{budkaEgyptianImpactPotBreaking2014}. These rituals, already present in Nubia since the New Kingdom, underscore a cultural continuity and exchange \parencites[p.~154, 165-166]{smithWretchedKushEthnic2003}[p.~120]{gasperiniAmaraWestCemeteries2023}.


After the Kushite kings withdrew from Egypt around 664 BCE, textual and archaeological sources confirm the ongoing presence of Nubian officials during the 26th Dynasty. Notable figures include the God's Wife of Amun, the High Priest of Amun and grandson of Shabaqo, the wife of the Mayor of Thebes, who was also the granddaughter of Piankhy; and Peksater, one of the wives of Piankhy. Furthermore, Egyptians began incorporating Nubian features into their burials during this time. These include the use of bead nets and offering tables, which were absent from Egyptian practices until their reintroduction by the Kushites \parencites[p.~514]{budkaKushiteTombGroups2010}[p.~494-498]{howleyKushitesEgypt6642020}.


During the 6th century BCE, relations between Nubia and Egypt were strained, likely due to the invasion of Lower Nubia by the 26th Dynasty king Psamtek II. During this time, the use of Egyptian writing became less common. Trade relations between Nubia and Egypt were reestablished around 570---526 BCE. Following the Persian conquest of Egypt, a text from the reign of Darius I (ca. 522---486 BCE) records a shipment of ivory from Kush. Additionally, it is likely that Nubians served in Xerxes I's army. During the reign of Harsiyotef (ca. 404---369 BCE)---one of the last kings of the Napatan period---, a stela inscribed with Egyptian hieroglyphics documents military campaigns in Lower Nubia. The text details battles against non-Egyptian groups, highlighting the region's distinctly intercultural nature even after its integration into Egypt \parencites[p.~344-345, 344-345, 376-377]{törökKingdomKushHandbook1997}[p.~361-370]{törökTwoWorldsFrontier2009}[p.~498-499]{howleyKushitesEgypt6642020}.


Beyond the royal pyramids, non-elite tombs from the later Napatan period are absent from the archaeological record. This absence can be attributed to several factors. Only Reisner's work at the Meroe cemeteries included non-royal burials dated to the Late Napatan phases, but these were exclusively elite populations. Therefore, the excavation methods of the time may have been inadequate to identify and date the simpler graves associated with non-elites. Additionally, the lack of datable material culture in low status tombs might have caused them to remain undated or even misdated to earlier phases. Other possibilities include the existence of undiscovered later Napatan cemeteries and the reuse of tombs over time, which further hinders their identification.


\section{Conclusion}


In this chapter, I explored the extensive history of cultural interactions between Nubia and Egypt, which shaped both cultures for millennia. This historical background is crucial as it demonstrates that the presence of Egyptian-style material culture in Napatan burials is not just an isolated occurrence, but rather a continuation of a long-standing process of cultural exchange. In particular, the position of Nubia in relation to Egypt during the New Kingdom period significantly influenced the development of material culture of the Nubian populations in subsequent years. This complex historical context is essential for its influence in material culture assemblages of the Napatan period, specially how different people selectively adopted and adapted certain cultural markers, reflecting their responses to shifting dynamics with Egypt. The historical background also reveals the intricate and evolving relationships between Kush and Egypt over time during the Napatan period. These ranged from cycles of Kushite control over Upper Egypt, periods of trade and conflict with the Assyrians, and to religious and political relations with the Amun cult in Thebes. These shifts set a stage for how amulet use might have varied across different social groups and chronological phases, depending on the political and cultural ties with Egypt. By examining these dynamics through a chronological and regional lens, we can better understand how amulet use shifted in response to these changing contacts with Egypt.
